\begin{definition}{}{}
Eine Kurve $\gamma : I \to \mathbb{R}^n$ ist eine stetige Abbildung, wobei $I \subset \mathbb{R}$ ein Intervall ist.\\
Wir schreiben


\[
\gamma(t) = \bigl( \gamma_1(t),\, \gamma_2(t),\, \dots,\, \gamma_n(t) \bigr),
\]


wobei jede Komponente $\gamma_i : I \to \mathbb{R}$ stetig ist.
\end{definition}

\begin{beispielbox}
\begin{enumerate}
  \item \textbf{Gerade:} Eine Gerade in $\mathbb{R}^n$ durch einen Punkt $a \in \mathbb{R}^n$ in Richtung eines Vektors $v \in \mathbb{R}^n$ wird beschrieben durch
  

\[
  \gamma(t) = a + t\,v, \quad t \in \mathbb{R}.
  \]


  
  \item \textbf{Kreis:} Ein Kreis in $\mathbb{R}^2$ mit Mittelpunkt $a \in \mathbb{R}^2$ und Radius $r > 0$ wird dargestellt durch
  

\[
  \gamma(t) = a + r\,(\cos t, \sin t), \quad t \in [0, 2\pi].
  \]


  
  \item \textbf{Helix:} Eine Helix in $\mathbb{R}^3$ kann beispielsweise definiert werden als
  

\[
  \gamma(t) = \bigl( r\cos t,\, r\sin t,\, c\,t \bigr), \quad t \in \mathbb{R},
  \]


  wobei $r > 0$ und $c \neq 0$ sind.
  
  \item \textbf{Graph einer Funktion:} Der Graph einer stetigen Funktion $f : I \to \mathbb{R}$ ist ebenfalls eine Kurve, nämlich
  

\[
  \gamma(t) = \bigl( t, f(t) \bigr), \quad t \in I.
  \]


\end{enumerate}
\end{beispielbox}